% ============================================================================
% Chapter 31 --- Institutional precedents
% ----------------------------------------------------------------------------
% Purpose : annotated catalogue of named institutions whose publicly
%           documented practices feed the architecture and the playbook.
%           Designed to be expanded by the international university
%           network in phase 2.
% Page target: ~5 pages.
% Dependencies: cited from Ch. 1 (synthesis), Ch. 18-27 (real-world
%               reference boxes).
% ============================================================================

\chapter{Institutional precedents}
\label{ch:institutional-precedents}

\begin{caveatbox}[Status of this chapter]
\noindent The institutional precedents in this chapter are
identified as starting points for the institution's own research,
not as direct templates. The architectural choices, operational
practices, and contractual arrangements at any cited institution
should be confirmed against that institution's own current public
documentation before being cited as precedent in derivative work.
This chapter records the institutions that have appeared in
public discussion of the patterns this Reference draws on; it
does not claim independent verification of every aspect of every
deployment.
\end{caveatbox}

The architecture and playbook of this Reference do not invent
their structural patterns from first principles --- they
consolidate patterns visible at research institutions whose practices
have been the subject of public discussion. This chapter records
the institutional patterns that have appeared most consistently as
references, organised by region and by the practice they
illustrate. The catalogue is intended for the international
university network anticipated in phase~2 of this initiative as a
starting point that the network will refine and extend with
direct contributions from each member.

%-----------------------------------------------------------------------------
\section{Method}
\label{sec:ip-method}

The institutions below are identified through three methodological
filters: appearance in research-and-education-network community
fora and at regional higher-education associations (such as EDUCAUSE
in North America) in recent publications; reference in NIST or ISO/IEC publications
that draw on academic deployments;
mention in publicly available research-network case material
(regional and national research-network case studies). The filters
bias the catalogue toward institutions that
publish their architectures --- in practice, larger and flagship
institutions that maintain public architecture documentation;
institutions whose practice is mature but undocumented, and smaller or
teaching-led institutions that consume the same primitives as shared
services rather than documenting their own builds, are necessarily
under-represented. This is a documentation bias, not a statement about
which institutions the architecture can serve (see
\S\ref{sec:ip-beyond-flagship}).

\paragraph{What is and is not claimed.} The catalogue claims that
each pattern has been the subject of public discussion at the kind
of institution to which it is attributed. The catalogue does \emph{not} claim
that the institution's current architecture matches the description;
that the institution endorses the pattern's use in this Reference;
or that the cited practice is the institution's authoritative
position. Each entry should be treated as a starting point for the
reader's own research.

\paragraph{Refresh cadence.} The catalogue would benefit from
annual review: institutional architectures evolve, public
documentation is revised or withdrawn, NREN community discussions
move forward. The phase-2 international network is the natural
locus for this refresh.

%-----------------------------------------------------------------------------
\section{European institutions}
\label{sec:ip-europe}

\paragraph{A large research organisation.} A large European
research organisation operating a multi-institutional
collaboration across many countries illustrates the pattern of an
identity federation, an audit pipeline, and a high-performance
experimental network architecture maintained under a strong
public-documentation practice. The migration of such an
organisation's identity and access management from a proprietary
federation service to a Keycloak-based open-source single
sign-on is documented in conference
proceedings~\cite{cern-keycloak-2020}. Scale and an open
documentation practice make this class of organisation one of the
most frequently cited research-organisation references for
research-intensive university IT.

\paragraph{A federal institute of technology.} A federal
institute of technology illustrates the model of granting local
administrator rights through a documented procedure (request, IT
risk assessment, approval by an authorised person in the
applicant's organisational unit) and an IT-services group
federated across departments. The pattern of distributed IT
authority within a coherent institutional framework is a
recurring point of reference in this Reference's
federated-governance discussions.

\paragraph{A collegiate university.} At a collegiate university,
a central information-services function operates the institution's
technology while an institutional data-network backbone connects
the central services and the constituent colleges. The collegiate
model illustrates a federated-network pattern under a coherent
institutional architecture.

\paragraph{A second collegiate university.} Distributed IT under
a central information-security policy framework is similarly
documented at other collegiate institutions. The pattern is
similar in spirit to the preceding case but differs in
operational details that each institution's own documentation
should be consulted for.

\paragraph{A national flagship research university.} A national
flagship research-intensive university operating a central IT
function is a typical European reference institution for
research-data infrastructure. Public documentation at such an
institution is often more limited than at the larger
network-anchored institutions.

%-----------------------------------------------------------------------------
\section{North American institutions}
\label{sec:ip-north-america}

\paragraph{A research-intensive institute of technology.} At a
research-intensive institute of technology, a central
information-systems function, an information-technology governance
committee reporting to the chief academic and executive officers,
and a flagship laboratory's technical-infrastructure group appear
recurrently as references for the federated-IT-governance pattern.
Such a governance committee typically operates with a set of
related committees (infrastructure steering, policy, project) that
consult on different aspects of IT strategy. Public documentation
at this class of institution is extensive and accessible.

\paragraph{A large research-intensive university.} A
campus-network function, a central university-IT organisation, and
a dedicated research-computing centre appear as references for
large-scale research-intensive university IT. The public
documentation of the multi-zone segmentation pattern at such
institutions is a frequently cited starting point.

\paragraph{A network-anchored public research university.} A
central IT-services organisation, an information-technology
council, and a regional research-network relationship illustrate
the network-anchored pattern at a large research-intensive
university. A college of engineering's own computing network is a
typical instance of faculty-level IT autonomy under institutional
governance.

\paragraph{A national research-cloud deployment.} A
research-cloud architecture~\cite{jetstream2-2021}, operated under
a national science-funding agency's access programme, is one of
the most extensively documented research-cloud deployments. It is
cited recurrently in this Reference's discussion of compute and
research-storage patterns.

\paragraph{A federated research university.} A research-intensive
university is frequently cited for its identity-federation
participation and for the integration of faculty IT with central
services. Public documentation is accessible through such an
institution's IT-services pages.

%-----------------------------------------------------------------------------
\section{Asia-Pacific institutions}
\label{sec:ip-asia-pacific}

\paragraph{A regional federation hub.} A leading
national research university in the region has been cited for
federation patterns at the regional scale and for the integration
of multi-jurisdiction research collaborations. Public
documentation is often partial, with the accessibility of
English-language material varying across regional references.

\paragraph{A federation-participating national university.} A
national university's participation in its country's access
federation and in network-anchored research networking is a
publicly documented reference for Asia-Pacific research-network
patterns.

\paragraph{A research-intensive university with limited public
documentation.} Some of the most-cited references in the region
maintain only limited English-language public documentation
relative to the European and North American institutions
discussed above. Direct engagement with such an institution would
be necessary for any precedent claim.

%-----------------------------------------------------------------------------
\section{Latin American institutions}
\label{sec:ip-latin-america}

\paragraph{A national flagship research university.} A national
flagship research university is one of the principal references
for Latin American research-network patterns. Its participation in
the national research and education network and in the national
identity federation provides documented federation precedent.

\paragraph{A large public research university.} A national
flagship public research university is a reference for large-scale
public-university IT. Public documentation in the national
language is often extensive, with English-language material more
limited.

%-----------------------------------------------------------------------------
\section{NRENs and consortium initiatives}
\label{sec:ip-nrens}

National Research and Education Networks and their consortia are
the principal multipliers of the patterns this Reference
develops. They produce shared documentation, operate shared
services, and host community fora in which architectural
discussions take place.

\paragraph{The regional research-network backbone.} The
research-network backbone serving the institution's region
connects the national research and education networks of that
region and operates the federated-identity and federated-access
services on which the Reference's identity and attachment
patterns rest --- typically a federated single-sign-on
interconnect, a collaborative-identity service, and a roaming
network-access service. Its community fora are a principal venue
for the discussions feeding this Reference.

\paragraph{The institution's national research and education
network.} The national research and education network connects
the country's research and education community to the regional
backbone and operates institutional services, commonly including
a national computer security incident response capability. The
national network's role as a shared operator is relevant to
several of the architecture's mutualisation patterns.

\paragraph{Regional research-network consortia.} Beyond the
national network, regional research-network consortia operate
services analogous to the backbone's and serve as the principal
venues for community discussion of architectural patterns at
research-intensive universities in their respective regions.

\paragraph{Sector service organisations.} National and regional
research-and-education service organisations operate the country
or region's academic backbone network and support a large body of
community guidance on which the Reference's patterns draw; several
also operate the national or regional identity federation.

%-----------------------------------------------------------------------------
\section{The precedents illustrate the architecture; they do not gate it}
\label{sec:ip-beyond-flagship}

The catalogue above is dominated by large, research-intensive
institutions for the methodological reason stated in
\S\ref{sec:ip-method}: such institutions maintain extensive public
architecture documentation, so their practice is citable. This must
not be read as an admissibility bar. As established in
\trchap{ch:architectural-principles}{} (\S\ref{sec:ap-necessity}), the
functions the architecture discharges are a minimal necessary set for
any institution offering a digital infrastructure, several of them
legally compelled regardless of size; the precedents illustrate
\emph{how} those functions are realised at scale, they do not define
\emph{which} institutions the architecture serves.

\paragraph{The shared-service channel.} A national
research and education network exists to serve its country's entire
research and education community, not its flagship institutions alone;
national networks are mandated accordingly, a mandate documented for
European national operators, e.g.~\cite{restena,dfn,switch-nren},
with comparable mandates in other regions. Through such
operators, federated identity through the national or regional
identity-federation profile~\cite{edugain}, collaborative
identity~\cite{eduteams}, and federated network access are available
to a small or teaching-led institution as managed services, without
that institution operating the underlying infrastructure;
\loc{research-backbone-scale} carries the same primitives across
borders (in the EU, GÉANT~\cite{geant}). A precedent need not be a flagship to be
relevant: an institution that consumes federated identity through its
national research and education network is exercising exactly the
identity-layer pattern of \trchap{ch:wave1-identity}{}, at a
realisation it can sustain.

\paragraph{Implication for the catalogue.} The phase-2 international
network (\S\ref{sec:ip-method}) is the natural locus for recording
such adoptions explicitly and broadening the catalogue beyond the
publication-rich flagships that necessarily dominate it today;
sector-guidance bodies such as Jisc in the UK~\cite{jisc-guidance},
and their regional analogues elsewhere, already document these patterns
for institutions across the size range.

%-----------------------------------------------------------------------------
\section{Documented gaps}
\label{sec:ip-gaps}

Several areas of the architecture remain thinly documented in
public institutional sources:

\begin{itemize}[leftmargin=1.5em,nosep]
  \item \textbf{Policy plane deployments at academic scale.} The
    industrial OPA and Cedar deployments provide most public
    documentation; academic deployments are reported informally
    in NREN fora but are rarely the subject of detailed public
    architecture documents.
  \item \textbf{TAC-aware audit pipelines.} The canonical TAC
    schema and the cross-layer correlation pattern are not, at
    the time of writing, the subject of any single publicly
    documented academic deployment. Various aspects are visible
    at various institutions; the consolidated pattern is novel
    to this Reference.
  \item \textbf{Sovereignty grid practice.} The grid is novel to
    this Reference. Adjacent practice exists (procurement
    sovereignty clauses, exit-data-export contractual language)
    but a structured five-dimension grid applied across an
    institutional IT portfolio has not, to the author's
    knowledge, been the subject of public institutional
    documentation.
\end{itemize}

\noindent The phase-2 international university network is the
locus where these gaps would be addressed: institutions that
adopt the architecture would publish their experience, and the
catalogue of precedents would acquire entries that this version
cannot yet record.
